Electricity prices provide a natural setting for statistical analysis because they combine strong temporal regularities, abrupt stress episodes, and cross-market dependence. In European day-ahead markets, prices are determined on an hourly basis for the following day, which yields a regular time series that is well suited to linear time-series tools while also exhibiting extreme events such as very high prices or negative prices.

In this report, we focus on the French and German day-ahead electricity markets over the year 2025. Our goal is not to apply several disconnected methods to the same dataset, but rather to answer a single question: when market conditions become stressed, do the French and German price series become more dependent, and do extreme events propagate over time and across markets?

To address this question, we use a common cleaned dataset extracted from the ENTSO-E Transparency Platform and organize the analysis in four steps. We first study the marginal temporal structure of the hourly series using linear time-series methods. We then analyze the tails of the distribution using extreme-value tools. Next, we investigate cross-market dependence through copulas. Finally, we model the temporal clustering of extreme events with Hawkes processes. This progression reflects the logic of the course: margins before dependence, thresholds motivated by extreme-value analysis, and point-process modeling for temporal clustering that copulas alone cannot capture.
